El proyecto est\'a organizado de manera simple y funcional. Los archivos principales son:

\begin{itemize}
  \item \texttt{src/creditmetrics\_project/data.py}: concentra los insumos del ejercicio.
  \item \texttt{src/creditmetrics\_project/model.py}: contiene toda la l\'ogica cuantitativa y de exportaci\'on.
  \item \texttt{src/creditmetrics\_project/main.py}: expone la interfaz de l\'inea de comandos.
\end{itemize}

\section*{Qu\'e hace \texttt{data.py}}

En este archivo se definen:

\begin{itemize}
  \item el espacio de 18 estados de rating;
  \item la curva libre de riesgo del ejercicio;
  \item los spreads de cr\'edito por rating;
  \item la matriz de transici\'on anual;
  \item la especificaci\'on de los tres bonos del portafolio.
\end{itemize}

Es, en la pr\'actica, la capa de entrada del modelo.

\section*{Qu\'e hace \texttt{model.py}}

Este archivo implementa casi todo el trabajo cuantitativo. Sus bloques m\'as importantes son:

\begin{itemize}
  \item \texttt{price\_bond\_today}: calcula el precio actual de un bono con la curva libre de riesgo m\'as el spread correspondiente a su rating.
  \item \texttt{price\_bond\_at\_horizon\_given\_state}: calcula el valor del bono a 1 a\~no, condicionado al estado migrado.
  \item \texttt{normalized\_transition\_row}: ajusta la fila de probabilidades seg\'un la pol\'itica aplicada a \texttt{NR}.
  \item \texttt{single\_bond\_distribution}: genera la distribuci\'on discreta de cada bono.
  \item \texttt{compress\_distribution\_to\_cents}: agrupa valores a nivel centavos para poder convolucionar una PMF discreta limpia.
  \item \texttt{convolve\_pmfs}: suma distribuciones discretas independientes.
  \item \texttt{portfolio\_distribution\_bruteforce}: enumera los \(5{,}832\) escenarios posibles y sirve de validaci\'on exacta.
  \item \texttt{write\_outputs}: exporta CSV, JSON, Excel, gr\'aficas e informe visual.
\end{itemize}

\section*{Qu\'e hace \texttt{main.py}}

El archivo principal recibe los argumentos de ejecuci\'on, lanza \texttt{run\_analysis} y muestra en consola un resumen del portafolio y de cada bono.

\section*{Flujo operativo}

El flujo completo del programa es el siguiente:

\begin{enumerate}
  \item cargar datos del ejercicio;
  \item valuar el precio actual de cada bono;
  \item construir la distribuci\'on de valor a 1 a\~no por bono;
  \item resumir media, volatilidad y VaR individual;
  \item convolucionar las tres distribuciones;
  \item construir la PMF del portafolio por \textit{brute force};
  \item comparar ambos resultados;
  \item exportar tablas y reportes.
\end{enumerate}

La decisi\'on de separar datos, l\'ogica y CLI vuelve al repositorio f\'acil de modificar. Por ejemplo, si cambia la pol\'itica \texttt{NR}, la curva libre de riesgo o el conjunto de bonos, no hace falta reescribir el flujo general.
